\metatitle{Lorentzian polynomials}

\metadescription{Lorentzian polynomials, their properties and examples including Schur polynomials, chromatic symmetric polynomials, and connections to matroids.}


\section[lorentzianPolynomial]{Lorentzian polynomials}

\name{Petter Brändén} and \name{June Huh} introduced the notion of \defin{Lorentzian polynomials}.
All \hyperref[stablePolynomial]{stable polynomials} are Lorentzian, but the converse is not true for polynomials
of degree three or more.
For an introduction to the topic, see \href{https://www2.math.upenn.edu/~jhaglund/IPAC/2022.html}{these lectures by Brändén}.


\begin{definition}[Lorentzian polynomial]

A homogeneous polynomial $h(x_1,\dotsc,x_n)$ of degree $d$ is \defin{strictly Lorentzian} if
\begin{itemize}
 \item All coefficients of $h$ are positive, and
 \item The signature of (the degree $2$ polynomial)
 \[
   \frac{\partial}{\partial x_{i_1}}\frac{\partial}{\partial x_{i_2}} \dotsm \frac{\partial}{\partial x_{i_e}} h(x_1,\dotsc,x_n)
 \]
 is $(+,-,-,\dotsc,-)$ for any $i_1,\dotsc,i_e \in [n]$ where $e = d-2$.

 To compute the signature of a homogeneous polynomial $P(x_1,\dotsc,x_n)$ of degree 2,
 consider the symmetric matrix
 \[
  \left(  \frac{\partial}{\partial x_{i}}\frac{\partial}{\partial x_{j}} P \right)_{ij} \in \setR^{n \times n}.
 \]
 If it has only one positive eigenvalue, while the remaining are negative, then it has signature $(+,-,-,\dotsc,-)$.

\end{itemize}

A polynomial is called \defin{Lorentzian} if it is the limit of strictly Lorentzian polynomials.
There are two other equivalent conditions listed in \cite{HuhMatherneMezarosDizier2022}.

An equivalent definition for Lorentzian is:

\begin{itemize}
 \item All coefficients are non-negative.
 \item The support of the polynomial forms an \hyperref[MConvexity]{M-convex set}.
 \item All quadratic forms (as above) has at most one positive eigenvalue.
\end{itemize}

\end{definition}

The \hyperref[basisGeneratingPolynomial]{basis-generating polynomial} of a
\hyperref[matroid]{matroid} is Lorentzian.  More generally, integral
\hyperref[polymatroids]{polymatroids} give Lorentzian normalized support
polynomials.  For the support condition, saturated Newton polytopes, and
discrete convexity background, see
\hyperref[newtonPolytopes]{Newton polytopes and $M$-convexity}.

Determining if polynomials are Lorentzian can be done efficiently, see \cite{Chin2024x}.
This is in contrast with checking for \hyperref[stablePolynomial]{real stability} which is coNP-complete, see \cite{Chin2024x}.


The following operations preserve the Lorentzian property:
\begin{itemize}
\item Taking the multi-affine part of a Lorentzian polynomial.
\item Products of Lorentzian polynomials are Lorentzian.
\item Specializing variables in a multiaffine Lorentzian polynomial.
\end{itemize}



In \cite{HuhMatherneMezarosDizier2022}, several generalizations of \hyperref[schurS]{Schur polynomials}
are either proved or conjectured to be Lorentzian.
\begin{theorem}[See \cite{HuhMatherneMezarosDizier2022}]
For any integer partition $\lambda$, the normalized \hyperref[schurS]{Schur polynomial}
\[
   \sum_{ \mu } K_{\lambda \mu} \frac{\monomial_\mu}{\mu_1! \mu_2 ! \dotsm \mu_\ell !}
\]
is Lorentzian. Here, $K_{\lambda \mu}$ are the
\hyperref[kostkaCoefficient]{Kostka coefficients}.
\end{theorem}
Another proof of this is given in \cite{RossToma2023}.

\name{Apoorva Khare}, \name{Jacob P. Matherne}, and
\name[Avery St. Dizier]{Avery St.~Dizier} extend this representation-theoretic
picture to shifted characters of parabolic Verma modules over
$\mathfrak{sl}_{n+1}(\setC)$ \cite{KhareMatherneDizier2025x}.  More generally,
they show that shifted and normalized characters of the Kostant partition
function of any loopless multigraph on $[n+1]$ are Lorentzian, and they give
limits on how far such log-concavity phenomena can be extended.


It is conjectured that the corresponding statement holds also for skew Schur polynomials, \hyperref[key]{key polynomials}
and \hyperref[schubert]{Schubert polynomials}.
In \cite{ChenFanYe2024x}, 0-1-\hyperref[grothendieck]{Grothendieck polynomials}
are shown to be Lorentzian.


In \cite{ChinQin2025x}, the authors study the set of symmetric polynomials which are also Lorentzian.
This set is homeomorphic to a closed Euclidean ball.  They also give a
reduction scheme for testing Lorentzianity of symmetric polynomials,
semialgebraic descriptions in small degrees, and examples showing that some
natural symmetric operators do not preserve the Lorentzian property.



\begin{example}[Chromatic symmetric polynomials]

For \emph{abelian Dyck paths}, the \hyperref[chromaticQuasisymmetricHistory]{chromatic symmetric polynomial}
(in any number of variables) is Lorentzian, see \cite{MatherneMoralesSelover2022x}.

In contrast, there is a (non-abelian) unit interval graph for which
the chromatic symmetric polynomial is not Lorentzian, or normalized Lorentzian.
The area sequence for when this fails is $01121121$ --- this example was found by
\name[R.I. Liu]{Ricky Liu} and \name{Cynthia Vinzant}.
\end{example}


In \cite{RossSussWannerer2023x}, the authors describe a way to construct operators preserving
the Lorentzian property.

See \cite{BrandenLeake2023x,Ross2023x} for Lorentzian polynomials on fans.
See \cite{HuXiao2023x} for Lorentzian polynomials and intersection theory.


For Lorentzian polynomials from independence sets in graphs,
see \cite{BendjeddouHardiman2025}.
\name{Joonkyung Lee}, \name{Jaeseong Oh}, and \name{Jaehyeon Seo} use
Lorentzian polynomial methods to prove homomorphism inequalities for
antiferromagnetic graphs \cite{LeeOhSeo2025x}.  This includes progress on
coloring and independent-set extremal problems for regular graphs.
